% ============================================================================
% SEC02.TEX - Section 8.2: Membuat HUD Skor
% ============================================================================

\section{Membuat HUD Skor}

Kita ingin pemain melihat skor mereka bertambah saat menghancurkan asteroid.

\subsection{Setup Text UI}
\begin{enumerate}
    \item Di dalam Canvas, klik kanan > \textbf{UI} > \textbf{Text - TextMeshPro}. (Jika diminta import TMP Essentials, klik Import).
    \item Beri nama \texttt{ScoreText}.
    \item Posisikan di pojok kiri atas (gunakan Anchor Presets: Top-Left).
    \item Ubah tulisan menjadi "Score: 0".
\end{enumerate}

\subsection{Script GameManager}
Kita butuh pengelola skor terpusat. Buat GameObject kosong bernama \texttt{GameManager}, lalu buat script \texttt{GameManager.cs}:

\begin{lstlisting}[language={[Sharp]C}]
using UnityEngine;
using TMPro; // Wajib ada untuk Text Mesh Pro

public class GameManager : MonoBehaviour
{
    public TextMeshProUGUI scoreText;
    private int score = 0;

    // Pola Singleton Sederhana agar mudah diakses
    public static GameManager Instance;

    void Awake()
    {
        Instance = this;
    }

    public void AddScore(int amount)
    {
        score += amount;
        scoreText.text = "Score: " + score;
    }
}
\end{lstlisting}

\begin{tips}
Jangan lupa drag object \texttt{ScoreText} ke slot script GameManager di Inspector.
\end{tips}

\subsection{Update AsteroidController}
Ubah agar saat asteroid hancur, skor bertambah:
\begin{lstlisting}[language={[Sharp]C}]
// Di dalam OnTriggerEnter2D AsteroidController
if (other.CompareTag("Bullet"))
{
    // Tambah Skor 10 poin
    GameManager.Instance.AddScore(10); 
    
    // ... ledakan dan destroy ...
}
\end{lstlisting}
